\section{Sintesi proteica}

La sintesi proteica avviene in 2 passaggi fondamentali:
\begin{itemize}
  \item \textbf{trascrizione}: nel nucleo;
  \item \textbf{traduzione}: nel citoplasma.
\end{itemize}

% ---------------------------------------------------------------------------
\subsection{Trascrizione}

Sintesi di RNA a partire da uno stampo di DNA, ad opera dell'\textbf{RNA polimerasi}. Di un gene viene usato come stampo un solo filamento di DNA (\textbf{filamento stampo}); l'RNA sintetizzato ha sequenza complementare al filamento stampo e corrispondente (a parte T $\rightarrow$ U) al \textbf{filamento non stampo} (o filamento di RNA senso).

La sintesi avviene in un segmento di DNA localmente svolto, la \textbf{bolla di trascrizione}: pochi nucleotidi del filamento stampo si appaiano con l'estremità della catena di RNA nascente; l'RNA polimerasi catalizza lo svolgimento e il successivo riavvolgimento dell'elica di DNA man mano che avanza.

\rubrica{Le 3 fasi della trascrizione}
\begin{enumerate}
  \item inizio;
  \item allungamento;
  \item terminazione.
\end{enumerate}

\rubrica{Inizio}
Nel promotore tipico di \textit{E.\ coli} si distinguono una sequenza $-35$ e una sequenza $-10$ (ricca in AT). L'RNA polimerasi si lega alla sequenza $-35$ e inizia lo svolgimento dei filamenti di DNA nella sequenza $-10$; la trascrizione comincia all'interno della bolla di trascrizione, in un sito da cinque a nove paia di basi oltre la sequenza $-10$.

\rubrica{Allungamento}
L'RNA polimerasi si lega al DNA ed estende covalentemente la catena di RNA, spostandosi progressivamente a valle: mantiene una piccola regione di ibrido DNA/RNA, con un sito di svolgimento a valle e un sito di riavvolgimento a monte della bolla di trascrizione.

\rubrica{Terminazione (terminatore rho-indipendente)}
Il terminatore contiene una regione ricca in GC seguita da almeno sei coppie di basi AT. La trascrizione di questa sequenza genera una catena di RNA con segmenti ricchi in GC che si appaiano tra loro subito dopo la sintesi, formando una struttura a forcina. Questa struttura ritarda il movimento dell'RNA polimerasi lungo il DNA, determinando la terminazione della trascrizione nel tratto adiacente di AT e il rilascio della catena nascente di RNA.

% ---------------------------------------------------------------------------
\subsection{Maturazione del pre-mRNA}

\rubrica{Aggiunta del cap}
In uno stadio precoce della trascrizione di un gene ad opera dell'RNA polimerasi II, un \textbf{cap} di 7-metil guanosina (7-MG) viene aggiunto all'estremità 5' del pre-mRNA, subito dopo l'inizio del processo di allungamento, attraverso una successione di reazioni enzimatiche: guanilil transferasi (a partire da GTP), guanina-7-metil transferasi e, eventualmente, 2'-O-metil transferasi.

\rubrica{Aggiunta del poli(A)}
Le estremità 3' dei trascritti sono prodotte per taglio endonucleolitico a valle di un segnale di poliadenilazione (sequenza consenso AAUAAA, seguita da una sequenza ricca in GU). La \textbf{poli(A) polimerasi} aggiunge quindi la coda di poli(A) alla nuova estremità 3' così generata.

% ---------------------------------------------------------------------------
\subsection{Gli attori della traduzione}

\rubrica{Ribosomi e RNA ribosomiale}
I ribosomi sono un complesso ribonucleoproteico \textbf{80S}, formato da una subunità maggiore \textbf{60S} e una subunità minore \textbf{40S}.
\begin{itemize}
  \item l'rRNA costituisce circa l'80\% di tutti gli RNA presenti nella cellula ed è sintetizzato nel nucleolo;
  \item la sintesi dell'RNA 45S, seguita da modificazioni chimiche e associazione con proteine specifiche, porta alla formazione del complesso ribonucleoproteico 80S;
  \item l'80S va incontro a maturazione, formando un complesso 40S e un complesso 60S;
  \item il complesso 40S passa nel citoplasma, dove forma la subunità ribosomale minore;
  \item il complesso 60S va incontro a ulteriori processamenti, formando l'rRNA 28S e l'rRNA 5,8S; questi, insieme all'rRNA 5S (di origine extranucleolare) e ad altre proteine, si trasferiscono nel citoplasma e formano la subunità ribosomale maggiore.
\end{itemize}
Sulla subunità maggiore si trovano 3 siti: \textbf{E}, \textbf{P}, \textbf{A}. Il solco tra le due subunità contiene il sito di legame per l'mRNA; ciascuna subunità contribuisce ai siti di legame per le molecole di tRNA.

\rubrica{RNA messaggero (mRNA)}
Trasferisce l'informazione genetica dal DNA ai ribosomi presenti nel citoplasma. Ha vita media breve: dopo il processo di sintesi proteica va incontro a rapida degradazione. Un \textbf{poliribosoma} (o \textit{polisoma}) è un mRNA associato a più ribosomi contemporaneamente.

Struttura dell'mRNA maturo, dall'estremità 5' alla 3': cap, \textbf{5' UTR} (regione non tradotta, sequenze a monte del sito di inizio), \textbf{CDS} (sequenza codificante, dal codone di start al codone di stop), \textbf{3' UTR}, coda di \textbf{poli-A}.

\rubrica{RNA transfer (tRNA)}
Localizzato nel citoplasma. Consiste in una catena polinucleotidica ripiegata su sé stessa, con una caratteristica struttura a trifoglio.
\begin{itemize}
  \item l'estremità 3' sopravanza quella 5' di tre nucleotidi, sempre \textbf{CCA}, che costituiscono il sito accettore dell'aminoacido;
  \item dalla parte opposta si trova l'\textbf{anticodone}: una sequenza di tre basi specifica per un determinato aminoacido.
\end{itemize}

\rubrica{Aminoacil-tRNA sintetasi}
Gli aminoacidi vengono legati covalentemente alle rispettive molecole di tRNA dalle \textbf{aminoacil-tRNA sintetasi}, che utilizzano ATP come fonte di energia, formando l'\textbf{aminoacil-tRNA}.

% ---------------------------------------------------------------------------
\subsection{Traduzione}

\rubrica{Inizio}
Nei batteri, la fase di inizio coinvolge un ribosoma, un mRNA, un tRNA iniziatore (a cui è legata una formil-metionina, fMet) e diversi fattori di inizio proteici.
\begin{enumerate}
  \item la subunità ribosomale minore si lega all'mRNA in corrispondenza del codone di inizio AUG; la sequenza leader a monte dell'AUG aiuta il ribosoma a identificare la sequenza;
  \item il tRNA iniziatore si lega al codone di inizio e uno dei fattori di inizio viene rilasciato;
  \item quando la subunità ribosomale maggiore si lega alla subunità minore e i rimanenti fattori di inizio sono rilasciati, il complesso di inizio è completo.
\end{enumerate}

\rubrica{Allungamento}
Ad ogni ripetizione del ciclo viene aggiunto un aminoacido alla catena polipeptidica in crescita.
\begin{enumerate}
  \item la catena polipeptidica è legata covalentemente al tRNA che trasporta l'ultimo aminoacido inserito nella catena; questo tRNA occupa il sito P del ribosoma;
  \item un aminoacil-tRNA si lega al sito A mediante l'appaiamento complementare tra le basi dell'anticodone sul tRNA e quelle del codone sull'mRNA;
  \item la catena polipeptidica in crescita si stacca dalla molecola di tRNA che occupa il sito P e si unisce con un legame peptidico all'aminoacido legato al tRNA nel sito A;
  \item nel passaggio di \textbf{traslocazione}, il ribosoma si sposta di un codone verso l'estremità 3' dell'mRNA: la catena polipeptidica in formazione viene trasferita sul sito P, e il tRNA non più carico lascia il ribosoma dal sito E.
\end{enumerate}

\rubrica{Terminazione}
\begin{itemize}
  \item un segnale di stop (codone di stop: UAA, UAG o UGA) determina la terminazione della sintesi del polipeptide;
  \item quando il ribosoma arriva su un codone di stop, un fattore di rilascio proteico si lega al sito A;
  \item il fattore di rilascio idrolizza il legame tra la catena polipeptidica e il tRNA, determinando il rilascio del polipeptide dalla molecola di tRNA nel sito P;
  \item le parti rimanenti del complesso di traduzione si dissociano.
\end{itemize}

% ---------------------------------------------------------------------------
\subsection{Codice genetico}

\rubrica{Proprietà del codice genetico}
\begin{itemize}
  \item \textbf{degenerato} o \textbf{ridondante}: più codoni possono codificare per lo stesso aminoacido;
  \item \textbf{univoco}: ogni codone è specifico per un solo aminoacido;
  \item \textbf{universale}: tutti gli organismi viventi hanno lo stesso codice genetico (eccezioni: mitocondri, funghi e protisti).
\end{itemize}
% ---------------------------------------------------------------------------
\subsection{Regolazione dell'espressione genica}

I meccanismi di controllo della trascrizione possono essere \textbf{inducibili} o \textbf{reprimibili}.

\rubrica{Regolazione negativa}
Lo stato basale di trascrizione è ``acceso'' a meno che un repressore non lo ``spenga''.
\begin{itemize}
  \item \textbf{trascrizione inducibile}: il repressore è una proteina il cui legame al DNA è inattivato da un induttore; in assenza di induttore il repressore è attivo e blocca la trascrizione, in sua presenza il repressore si inattiva e la trascrizione avviene;
  \item \textbf{trascrizione reprimibile}: il repressore è formato dall'interazione tra una proteina \textbf{aporepressore} e un \textbf{corepressore}; in assenza di corepressore l'aporepressore è inattivo e la trascrizione avviene, in sua presenza il repressore si attiva e blocca la trascrizione.
\end{itemize}

\rubrica{Operone lac (controllo negativo inducibile)}
Nei batteri, il lattosio presente nella cellula si lega al repressore rendendolo inattivo e quindi \textbf{induce} la trascrizione dei geni che codificano per gli enzimi che idrolizzano il lattosio.

\rubrica{Operone trp (controllo negativo reprimibile)}
Nei batteri, il triptofano presente nella cellula si lega al repressore rendendolo attivo e quindi \textbf{reprime} la trascrizione dei geni che codificano per gli enzimi che sintetizzano il triptofano.

\rubrica{Regolazione positiva}
Lo stato basale della trascrizione è ``spento''. La trascrizione è stimolata dal legame di una proteina attivatrice trascrizionale al sito di legame dell'attivatore.

\rubrica{Esempio: regolazione positiva ad opera degli estrogeni}
\begin{itemize}
  \item gli estrogeni agiscono legandosi a recettori specifici che si trovano all'interno del nucleo;
  \item il complesso estrogeno-recettore, insieme a coattivatori, si lega a siti specifici del DNA (elementi di risposta agli estrogeni), stimolando la trascrizione di geni che codificano per specifiche proteine.
\end{itemize}

% ---------------------------------------------------------------------------
\subsection{Regolazione dell'espressione genica negli eucarioti}

Negli eucarioti la regolazione avviene su 4 livelli:
\begin{enumerate}
  \item \textbf{trascrizionale}: il rimodellamento della cromatina rende i geni accessibili per la trascrizione, e viene regolato l'inizio della trascrizione $\rightarrow$ determina quali geni vengono trascritti;
  \item \textbf{post-trascrizionale}: variazioni nella maturazione del pre-mRNA, rimozione delle proteine mascheranti, variazioni nella velocità di degradazione dell'mRNA, RNA interference $\rightarrow$ determina il tipo e la disponibilità degli mRNA per i ribosomi;
  \item \textbf{traduzionale}: variazioni nella velocità di inizio della sintesi proteica $\rightarrow$ determina la velocità alla quale le proteine vengono prodotte;
  \item \textbf{post-traduzionale}: variazioni nella velocità di maturazione delle proteine, rimozione dei segmenti mascheranti, variazioni nella velocità di degradazione delle proteine $\rightarrow$ determina la disponibilità delle proteine mature.
\end{enumerate}
